\input{../preamble.tex}
\SetDate[11/03/2026]
\setbeamertemplate{page number in head/foot}{}

\begin{document}

\section{Automatismes}

\begin{frame}

\centering \huge
Automatismes

\large
Calculatrice interdite

\end{frame}

\subsection{Fractions}

\framedelayed{
	Calculer. Exprimer sous forme de fraction irréductible.
	\boxAB{
		$\dfrac35 + \dfrac64$
	}{
		$\dfrac53 + \dfrac24$
	}
}{
	\boxAB{
		$\dfrac35 + \dfrac64 = \dfrac{3\times4 + 6\times5}{5\times4} = \dfrac{42}{20} = \dfrac{21}{10}$
	}{
		$\dfrac53 + \dfrac24 = \dfrac{5\times4+2\times3}{3\times4} = \dfrac{26}{12} = \dfrac{13}{6}$
	}
}

\subsection{Différence}

\framedelayed{
	Calculer la différence $x-y$ lorsque
	\boxAB{
		$x = -1$ et $y= -4$.
	}{
		$x = -2$ et $y = -7$.
	}
}{
	\boxAB{
		$-1 - (-4) = -1 + 4 = 3$
	}{
		$-2 - (-7) = -2 + 7 = 5$
	}
}

\subsection{Coordonnées de vecteur}


\framedelayed{
	\vspace{-12pt}
	\begin{center}
	\begin{tikzpicture}[>=stealth, scale=.83]
		\begin{axis}[xmin = -2.1, xmax=3.1, ymin=-2.1, ymax=3.1, axis x line=middle, axis y line=middle, axis line style=->, grid=both]
			\addplot[RED_E, very thick, ->] (1, 2) -- (axis cs:2,-1) node[midway, right, text=RED_E] {\Large $u$};
			\addplot[GREEN_E, very thick, ->] (1,1) -- (axis cs:-1,2)  node[midway, below, text=GREEN_E] {\Large $v$};
		\end{axis}
	\end{tikzpicture}
	\end{center}
	\vspace{-20pt}
	\boxAB{
		Donner approximativement les coordonnées de $u$.
	}{
		Donner approximativement les coordonnées de $v$.
	}
}{
	\begin{center}
	\begin{tikzpicture}[>=stealth, scale=.83]
		\begin{axis}[xmin = -2.1, xmax=3.1, ymin=-2.1, ymax=3.1, axis x line=middle, axis y line=middle, axis line style=->, grid=both, clip=false]
			\addplot[RED_E, very thick, ->] (1, 2) -- (axis cs:2,-1) node[pos=.3, right, text=RED_E] {\Large $u = \pvec1{-3}$};
			\addplot[GREEN_E, very thick, ->] (1,1) -- (axis cs:-1,2)  node[left, text=GREEN_E] {\Large $\pvec{-2}1 = v$};
		\end{axis}
	\end{tikzpicture}
	\end{center}
}

\subsection{Image}

\framedelayed{
	Soit $f(x) = (x+2)(x-1)$ une fonction sur $\R$.
	\boxAB{
		Calculer $f(2)$.
	}{
		Calculer $f(3)$.
	}
}{
	$f(x) = (x+2)(x-1)$
	\boxAB{
		$f(2) = (2+2)(2-1) = (4)(1) = 4$.
	}{
		$f(3) = (3+2)(3-1) = (5)(2) = 10$.
	}
}


\section{Solutions}

\shipoutAnswer

\end{document}